%قبل از شروع مسئله ترجیحاً چیزی ننویسید. اما اگر خواستید بنویسید، در قالب لیستی (فقرات و شمارش) بنویسید.
%برای اضافه کردن نمره، بعدی از شروع‌مسئله، [x] را اضافه کنید و به جای x نمره مورد نظر را قرار دهید. برای زیرمسئله‌ها هم به همین صورت است.
%\شروع‌مسئله[10]
\شروع‌مسئله 
در این مرحله شما یک الگوریتم \lr{space partitioning} را پیاده‌سازی می‌کنید. برای آشنایی اولیه با الگوریتم‌ها و داده‌ساختارهای \lr{space partitioning} می‌توانید از 
\href{https://en.wikipedia.org/wiki/Space_partitioning}{این لینک}
 و 
 \href{http://cs.ucf.edu/~dcm/Teaching/COT4810-Spring2011/Presentations/JonLeonardSpacePartitioningDataStructures.pdf}{این لینک}
  استفاده کنید. شما باید یکی از الگوریتم‌های \lr{space partitioning} را برای پیاده‌سازی انتخاب کنید و نام آن را در خط اول کدتان بنویسید. در 
  \href{https://en.wikipedia.org/wiki/List_of_data_structures}{این لینک}
   لیستی از داده‌ساختارهای \lr{space partitioning} آورده شده است. در این مرحله لیستی از مستطیل‌ها به شما داده می‌شود. شما باید تمام جفت مستطیل‌هایی که با هم برخورد دارند را خروجی دهید و در الگوریتمی که پیاده‌سازی می‌کنید از یک داده‌ساختار \lr{space partitioning} استفاده کنید. دقت کنید که مجاز به استفاده از کدهای موجود در اینترنت نیستید. تنها از زبان برنامه نویسی پایتون می‌توانید استفاده کنید و امکان استفاده از هیچ کتابخانه‌ای را ندارید. همچنین امکان استفاده از روش‌های \lr{multithreading} را ندارید و برنامه‌تان باید به صورت ترتیبی اجرا شود. به ازای یک ورودی برنامه‌تان باید  همیشه جواب ثابتی تولید کند. کد شما باید کاملاً درست کار کند. به این معنا که در صورتی که کدتان \lr{time limit} نخورد باید پاسخ درستی بدهد. پیشنهاد می‌شود که جست‌و‌جو در مورد الگوریتم‌های مختلف \lr{space partitioning} را به روزهای آخر موکول نکنید چرا که فرصت پیاده‌سازی، تست و بهبود آن‌ها را پیدا نخواهید کرد.
 \\
 \\
 \متن‌سیاه{ورودی و خروجی:}
در هر خط یک مستطیل به صورت $(x_{min}, x_{max}, y_{min}, y_{max})$ داده می‌شود و در آخر \lr{end} می‌آید. مختصات هر مستطیل در بازه‌ی $(-1000, 1000)$ است. در هر خط از خروجی یک برخورد به صورت $(index_1, index_2)$ چاپ می‌شود که $index$ها شماره‌های مستطیل‌ها هستند. شماره‌ی مستطیل‌ها به ترتیب ورودی دادن آن‌ها است و از $0$ شروع می‌شود. دقت کنید که مختصات مستطیل‌ها لزوما صحیح نیستند.
 \صفحه‌جدید
  نمونه:
  
  %\lstinputlisting{source_filename.py}
  \begin{latin}
  \begin{lstlisting}
  # input
  1 3 0 4
  5 9 4 7
  1 7 0 5
  2 12 2 8
  1 5 0 10
  1 3 0 4
  5 9 4 7
  end
  # output
  1 2
  1 3
  2 6
  4 5
  0 2
  0 3
  0 5
  2 3
  3 6
  0 4
  1 6
  2 5
  3 4
  2 4
  3 5
  \end{lstlisting}
  \end{latin}

%\متن‌سیاه{سوال عملی اول:}\\
%\متن‌سیاه{دراپ‌تست}\\

%\متن‌سیاه{ورودی:}\\
%در خط اول ورودی، ابتدا k و سپس m داده می‌شود.\\
%\متن‌سیاه{خروجی:}\\
%در تنها سطر خروجی، تعداد آیفون‌های مورد نیاز می‌آید. اگر در هیچ صورتی چنین کاری ممکن نبود ۱- چاپ می‌شود.\\
%\متن‌سیاه{محدودیت‌ها:}\\
%$1\leq k\leq 10^4$

%$m\leq10^9$
%\\
%\textbf{\متن‌سیاه{ورودی و خروجی نمونه}\\
%% تست اول
%\begin{latin}
% \begin{tabularx}{\textwidth}{|X | X|}
% \hline
% \textbf{stdin} &
% \textbf{stdout}\\
% \hline
% 100 100&
% 1\\
% \hline
% 
% \end{tabularx}
%\end{latin}
%
%% تست دوم
%
%\begin{latin}
% \begin{tabularx}{\textwidth}{|X | X|}
% \hline
% \textbf{stdin} &
% \textbf{stdout}\\
% \hline
% 340 20&
% 3\\
% \hline
% 
% \end{tabularx}
%\end{latin}}
% صورت اصلی مسئله را در این‌جا بنویسید.
%اگر مسئله چند قسمت داشت به این صورت عمل کنید:
%\شروع{زیرمسائل}
%\زیرمسئله[0] % مسئلۀ ۱آ
%\زیرمسئله % مسئلۀ ۱ب
%\زیرمسئله % مسئلۀ ۱پ
%\زیرمسئله % مسئلۀ ۱ت
%\پایان{زیرمسائل}




%%%%%%%%%%%%%%%%%%%%%%%%%%%%%%%%%%%%%%%%%%%%%%%%%%%%%
% برخی دستورات مفید برای تمرین‌ها. %
%%%%%%%%%%%%%%%%%%%%%%%%%%%%%%%%%%%%%%%%%%%%%%%%%%%%%

% فرمول‌های نمونه:
%$f(x) = \Theta(x)$
%$T(x, y) \leq \log(x) + 2^y + \binom{2n}{n}$
% {\tt code\_function}


% لیست موارد بدون شماره‌گذاری:
%\شروع{فقرات}
% \فقره مورد اول 
% \شروع{فقرات}
% \فقره مورد اول داخلی 
% \شروع{فقرات}
% \فقره اولین مورد خیلی داخلی 
% \فقره دومین مورد خیلی داخلی 
% \پایان{فقرات}
% \فقره مورد دوم داخلی 
% \پایان{فقرات}
% \فقره مورد دوم 
%\پایان{فقرات}

% لیست موارد شماره‌گذاری شده:
%\شروع{شمارش}
% \فقره یک مورد 
% \فقره یک مورد دیگر 
% \فقره باز هم یک مورد دیگر
%\پایان{شمارش}

% برای نوشتن یک فرمول چند خطی که یک شماره بخوره. برای حالت‌های دیگر به سایت زیر مراجعه کنید:
% https://www.overleaf.com/learn/latex/Aligning_equations_with_amsmath
%\begin{align} 
% \begin{split}
% (x+y)^3 &= (x+y)^2(x+y) \\
% &= (x^2+2xy+y^2)(x+y) \\
% &= (x^3+2x^2y+xy^2) + (x^2y+2xy^2+y^3) \\
% &= x^3+3x^2y+3xy^2+y^3
% \end{split} 
%\end{align}

% نوشتن ماتریس و آرایه دوبعدی
%\begin{align}
% A = 
% \begin{bmatrix}
% A_{11} & A_{21} \\
% A_{21} & A_{22}
% \end{bmatrix}
%\end{align}

% اضافه کردن تصویر با شرح. برای اضافه کردن تصویر بدون شرح، فقط دستور درج‌تصویر را استفاده کنید.
%\شروع{شکل}[h!]
%\تنظیم‌ازوسط
% \درج‌تصویر[پهنا=0.5\پهنای‌متن]{img.jpg}
%\شرح{شرح تصویر}
%\برچسب{تصویر:۱}
%\پایان{شکل}
%\پایان{زیرمسائل}

%کد زیر را ببینید.
%%\lstinputlisting{source_filename.py}
%\begin{latin}
%\begin{lstlisting}
%# Import the modules
%import sys
%import random
%
%ans = True
%
%while ans:
%question = raw_input("Ask the magic 8 ball a question: (press enter to quit) ")
%
%answers = random.randint(1,8)
%
%if question == "":
%sys.exit()
%
%elif answers == 1:
%print "It is certain"
%
%elif answers == 2:
%print "Outlook good"
%
%elif answers == 3:
%print "You may rely on it"
%
%elif answers == 4:
%print "Ask again later"
%
%elif answers == 5:
%print "Concentrate and ask again"
%
%elif answers == 6:
%print "Reply hazy, try again"
%
%elif answers == 7:
%print "My reply is no"
%
%elif answers == 8:
%print "My sources say no"
%\end{lstlisting}
%\end{latin}